% SPDX-License-Identifier: LicenseRef-ECHO-NC-Research-1.0
% Options for packages loaded elsewhere
\PassOptionsToPackage{unicode}{hyperref}
\PassOptionsToPackage{hyphens}{url}
\documentclass[
]{article}
\usepackage{xcolor}
\usepackage{amsmath,amssymb}
\setcounter{secnumdepth}{-\maxdimen} % remove section numbering
\usepackage{iftex}
\ifPDFTeX
  \usepackage[T1]{fontenc}
  \usepackage[utf8]{inputenc}
  \usepackage{textcomp} % provide euro and other symbols
\else % if luatex or xetex
  \usepackage{unicode-math} % this also loads fontspec
  \defaultfontfeatures{Scale=MatchLowercase}
  \defaultfontfeatures[\rmfamily]{Ligatures=TeX,Scale=1}
\fi
\usepackage{lmodern}
\ifPDFTeX\else
  % xetex/luatex font selection
\fi
% Use upquote if available, for straight quotes in verbatim environments
\IfFileExists{upquote.sty}{\usepackage{upquote}}{}
\IfFileExists{microtype.sty}{% use microtype if available
  \usepackage[]{microtype}
  \UseMicrotypeSet[protrusion]{basicmath} % disable protrusion for tt fonts
}{}
\makeatletter
\@ifundefined{KOMAClassName}{% if non-KOMA class
  \IfFileExists{parskip.sty}{%
    \usepackage{parskip}
  }{% else
    \setlength{\parindent}{0pt}
    \setlength{\parskip}{6pt plus 2pt minus 1pt}}
}{% if KOMA class
  \KOMAoptions{parskip=half}}
\makeatother
\usepackage{longtable,booktabs,array}
\usepackage{calc} % for calculating minipage widths
% Correct order of tables after \paragraph or \subparagraph
\usepackage{etoolbox}
\makeatletter
\patchcmd\longtable{\par}{\if@noskipsec\mbox{}\fi\par}{}{}
\makeatother
% Allow footnotes in longtable head/foot
\IfFileExists{footnotehyper.sty}{\usepackage{footnotehyper}}{\usepackage{footnote}}
\makesavenoteenv{longtable}
\setlength{\emergencystretch}{3em} % prevent overfull lines
\providecommand{\tightlist}{%
  \setlength{\itemsep}{0pt}\setlength{\parskip}{0pt}}
\usepackage{bookmark}
\IfFileExists{xurl.sty}{\usepackage{xurl}}{} % add URL line breaks if available
\urlstyle{same}
\hypersetup{
  pdftitle={Does a similarity recommender shape simulated taste? A pre-registered simulation with an assumption-sensitivity test and a real-catalogue second test},
  hidelinks,
  pdfcreator={LaTeX via pandoc}}

\title{Does a similarity recommender shape simulated taste? A
pre-registered simulation with an assumption-sensitivity test and a
real-catalogue second test}
\author{}
\date{}

\begin{document}
\maketitle

\textbf{Status: independent research manuscript / preprint draft, not
peer-reviewed.} Release candidate 0.1.0-rc1, 2026-11-26. All findings
are simulation results; simulated preference change is not evidence of
human persuasion.

\subsection{Author and affiliation
metadata}\label{author-and-affiliation-metadata}

Mohamed Bechir Kefi, independent researcher (owner-approved; owner,
2026-11-26, decision sheet). An AI system (Claude, Anthropic) assisted
with research engineering, analysis code and drafting as a disclosed
research and development assistant. It is not an author. See section 11.

Data: FMA metadata (Defferrard et al., arXiv:1612.01840), CC BY 4.0. No
audio was downloaded and no listener data exists in this study.

\subsection{Abstract}\label{abstract}

Recommender systems decide what listeners are exposed to; if tastes
adapt to what is consumed, the recommender may shape the preferences it
is meant to reflect. We ask, in a seeded simulation, whether a
personalised similarity policy (SIM) changes simulated listeners'
end-of-session taste differently from uniform-random exposure (RAND).
The protocol, decision rule and smallest effect size of interest (0.05
rad) were frozen by the owner before any confirmatory run. The primary
experiment (EXP-002; 5 policies, fixed and adaptive preferences, 10
seeds, synthetic catalogue calibrated on FMA metadata) found that SIM
produced less angular preference drift than RAND: Delta = -0.0355 rad,
95 \% bootstrap CI {[}-0.0535, -0.0200{]}, so H0 was rejected; whether
the difference is practically relevant is unresolved. The pre-registered
sensitivity experiment (EXP-003, 360 runs) withdrew this interpretation
as assumption-dependent: under a bounded-confidence update rule at the
same rate the contrast was practically equivalent to zero (+0.0069 rad,
CI {[}-0.0130, 0.0303{]}) or reversed with an oracle profile (+0.1883
rad, CI {[}0.1598, 0.2196{]}), and the sign also changed with the update
rate. A pre-specified second test on 4,951 real FMA tracks (EXP-004)
reversed the sign (+0.1652 rad, CI {[}0.1560, 0.1735{]}): the result
does not replicate. In descriptive secondary outcomes, SIM concentrated
exposure more than RAND on both catalogues, as a personalised top-k
policy does largely by construction. Within this model, the direction of
simulated taste change is not a property of the policy alone; it depends
on the update rule, the update rate and the catalogue geometry.

\subsection{Keywords}\label{keywords}

recommender systems; feedback loops; preference drift; simulation;
pre-registration; sensitivity analysis; music recommendation

\subsection{1. Introduction}\label{introduction}

A recommender that learns from consumption can also change what is
consumed. If listeners' tastes adapt to what they hear, a policy could
shape taste rather than simply reflect it. Simulation studies have
described homogenisation and preference amplification under such
feedback loops {[}chaney2018, mansoury2020, kalimeris2021, jiang2019{]},
and recent work defines a policy-induced preference shift relative to a
``natural'' reference {[}carroll2022{]}. Of the simulation studies read
for this work, two hold the users' intrinsic preferences fixed and
measure homogenisation or a shift in the system's representation of
taste {[}chaney2018, mansoury2020{]}, and two model one specified form
of preference or interest dynamics {[}kalimeris2021, jiang2019{]}. In
the limited, non-systematic set of sources read, we did not find a
pre-registered test of whether a policy comparison survives a change of
that mechanism; this is a statement about our reading, not a claim that
no such test exists.

\textbf{Research question.} Can recommendation policies change exposure
and simulated taste rather than simply reflect it? Operationally: in a
seeded simulation of listeners with fixed or adaptive preference
vectors, does the policy change end-of-session preferences beyond what a
non-personalised uniform-random policy induces, and does it change what
listeners are exposed to?

\textbf{Hypotheses} (frozen, protocol v1.0). Primary estimand: under
adaptive preferences, Delta = E\_seed{[}mean\_u theta\_u(SIM) - mean\_u
theta\_u(RAND){]}, in radians, where theta\_u is the angle between
listener u's initial and final preference vectors. H0: Delta = 0; H1:
Delta != 0 (non-directional, because the literature motivates both
narrowing and amplification).

\textbf{Contributions} (what was actually done):

\begin{enumerate}
\def\labelenumi{\arabic{enumi}.}
\tightlist
\item
  A small, fully reproducible, seed-paired comparison of five simple
  policy families against uniform-random exposure, with fixed and
  adaptive preferences modelled separately (EXP-002, 100 runs).
\item
  A pre-registered retain-or-withdraw sensitivity test of the primary
  contrast across three update rules, three update rates and two
  recommender profiles (EXP-003, 360 runs). It withdrew the primary
  interpretation.
\item
  A pre-specified second test of the primary contrast on a real music
  catalogue (FMA test split, real genres, audio features and play
  counts, simulated listeners; EXP-004, 100 runs). It did not replicate
  the primary direction.
\item
  An open evidence trail: frozen protocols with hashes, per-run
  manifests, byte-reproducible analysis outputs and a claims register.
\end{enumerate}

No novelty claim is made beyond this specific combination; the
literature search was targeted, not systematic (section 2).

\subsection{2. Related work}\label{related-work}

The sources below were opened and read in full (seven on 2026-11-25, at
M6; the FMA licence statement comes from the repository README); the
section or page that supports each statement is given in
\texttt{research/publication/claims\_audit.md} (M6 reference table). The
search was targeted and non-systematic
(\texttt{research/literature\_review.md}); citation chaining from
kalimeris2021 and carroll2022 was not completed.

\begin{itemize}
\tightlist
\item
  \textbf{Feedback loops and homogenisation.} Chaney, Stewart and
  Engelhardt {[}chaney2018{]} report a simulation in which algorithmic
  confounding in recommendation feedback loops homogenises user
  behaviour without raising utility; their users' preferences are fixed
  (their assumption 3). Mansoury et al.~{[}mansoury2020{]} simulate
  popularity-bias amplification with reduced aggregate diversity and a
  shift in the representation of user taste; they state that this shift
  is in the system's representation of preferences, not in intrinsic
  preferences (their section 4.3.2). ECHO reuses the idea of a
  homogenisation outcome and includes a popularity arm (POP) and an
  exposure-concentration metric.
\item
  \textbf{Preference dynamics.} Jiang et al.~{[}jiang2019{]} study
  degenerate feedback loops in user-interest dynamics and distinguish
  echo chambers from filter bubbles. Kalimeris et
  al.~{[}kalimeris2021{]} model preference amplification in a
  matrix-factorisation recommender, with users drifting toward
  recommended content. ECHO's consumed-item update rule is a simple
  member of this family. ECHO does not replicate either model; it uses a
  simpler angular-drift model.
\item
  \textbf{Reference for induced shift.} Carroll et al.~{[}carroll2022{]}
  evaluate induced preference shifts against ``safe'' shifts and
  operationalise ``natural preference shifts'' as the preferences a user
  would have when interacting with a random recommender, an
  approximation of having no recommender at all (their section 5).
  ECHO's uniform-random comparator (RAND) follows this
  operationalisation. As in their setting, RAND itself moves preferences
  under adaptive preferences (section 5.1), so the contrast measures
  shift beyond the random reference, not shift relative to no change.
\item
  \textbf{Simulation platforms.} RecSim {[}ie2019{]} is a configurable
  simulation platform whose user latent state and its dynamics are
  configurable. It served as a design reference only and is not a
  dependency.
\item
  \textbf{Metrics.} Vargas and Castells {[}vargas2011{]} (full text read
  on 2026-11-25) unify novelty and diversity metrics. Their equation 5
  defines item novelty as self-information, -log2 p(i \textbar{} seen),
  and their mean self-information and average intra-list distance
  (equation 16, without rank discount and relevance weighting) match
  ECHO's novelty (base-popularity share as p) and intra-list diversity
  (distance 1 - cos).
\item
  \textbf{Data.} FMA {[}defferrard2017{]} (paper and repository README
  read) provides the metadata, genres, audio features and play counts
  used for calibration and for EXP-004.
\end{itemize}

\textbf{Positioning.} Relative to these works, ECHO is an extension in
method, not in mechanism: it adds a frozen protocol, a negative control,
a pre-registered withdrawal rule and a real-catalogue second test to a
deliberately simple model. Its main finding is negative for the
generality of the primary contrast.

\subsection{3. Background and formal
model}\label{background-and-formal-model}

The full specification is \texttt{research/protocol.md} sections 3, 4
and 10; implementation choices frozen with the code are in DEC-006.

\begin{itemize}
\tightlist
\item
  \textbf{Catalogue.} N = 1000 tracks with unit-norm embeddings x\_i in
  R\^{}d, d = 16. In the synthetic catalogue (EXP-002, EXP-003),
  embeddings are drawn around K = 16 Gaussian cluster centres with
  cluster weights fitted on the FMA training split, then L2-normalised;
  base popularity follows a Zipf law with exponent s = 1.3065, fitted on
  the same split (journal J-0003). In EXP-004, each seed samples 1000 of
  4,951 eligible FMA test-split tracks; embeddings are the first 16
  principal components of standardised librosa statistics
  (training-split basis), L2-normalised; clusters are the 16 FMA
  top-level genres; popularity is listens + 1.
\item
  \textbf{Listeners.} U = 200 per seed. The initial preference p\_u,0 is
  a unit vector near a uniformly chosen cluster centre.
\item
  \textbf{Session.} T = 100 steps. Each step, the policy shows a slate
  of k = 10 tracks from the candidates the listener has not consumed
  (candidate exclusion). The listener picks one track or skips, by a
  softmax over cos(p\_u,t, x\_i)/tau with tau = 0.1 and a skip option of
  utility 0.
\item
  \textbf{Preference models.} Fixed: p\_u,t+1 = p\_u,t. Adaptive,
  consumed-item rule (primary): p\_u,t+1 = normalise((1 - eta) p\_u,t +
  eta x\_chosen), eta = 0.05, no update on skip. EXP-003 adds an
  exposure rule (move toward the slate mean) and a bounded-confidence
  rule (update only if cos(p\_u,t, x\_chosen) \textgreater{} 0.5).
\item
  \textbf{Policies.} RAND: uniform-random slate (baseline). POP: top-k
  by global consumption count plus base popularity. SIM: top-k by
  cos(p\_hat, x\_i), where p\_hat is the normalised mean of consumed
  embeddings after a 5-step random warm-up (the recommender never sees
  p\_u,t, except in the oracle ablation of EXP-003). MMR: maximal
  marginal relevance re-ranking of SIM, lambda = 0.5. EXPLORE:
  epsilon-greedy SIM, epsilon = 0.2.
\item
  \textbf{Primary metric.} theta\_u = arccos(clip(\textless p\_u,0,
  p\_u,T\textgreater, -1, 1)), computed stably as 2
  arcsin(\textbar\textbar p\_u,0 - p\_u,T\textbar\textbar{} / 2),
  averaged over listeners. Its range is {[}0, pi{]}; two unrelated
  directions in R\^{}16 lie close to pi/2 apart, so drift near pi/2
  means the final taste is nearly unrelated to the initial one.
\item
  \textbf{Secondary metrics.} Homogenisation (change in mean pairwise
  cosine of preferences), discovery (share of clusters consumed),
  novelty (mean self-information of consumed tracks, bits), intra-list
  diversity (mean pairwise 1 - cos within slates), relevance to initial
  and current taste, exposure Gini over slate impressions, skip rate,
  and drift per accepted item.
\item
  \textbf{Cost.} One run scores U x T x N = 2 x 10\^{}7 cosines; runs
  take seconds on a laptop CPU.
\end{itemize}

\subsection{4. Methods}\label{methods}

\begin{itemize}
\tightlist
\item
  \textbf{Design.} EXP-001 (exploratory): simulator verification and
  smoke matrix (checks V1 to V6: determinism, zero drift under fixed
  preferences, candidate exclusion, unit norms, metric bounds, schema).
  EXP-002 (confirmatory, primary): 5 policies x \{fixed, adaptive eta
  0.05\} x 10 seeds = 100 runs. EXP-003 (confirmatory sensitivity): RAND
  and SIM x 3 update rules x eta in \{0.01, 0.05, 0.2\} x \{estimated,
  oracle\} profile x 10 seeds = 360 runs. EXP-004 (confirmatory
  secondary, second test): the EXP-002 matrix on the FMA catalogue, 100
  runs.
\item
  \textbf{Unit and pairing.} The seed (one catalogue and one listener
  population) is the independent unit; listeners within a seed share the
  catalogue and POP counts, so they are aggregated to a seed mean.
  Common random numbers: every policy within a seed sees the same
  catalogue, listeners and choice uniforms. Confirmatory seeds 20261124
  to 20261133 were fixed in advance and disjoint from exploratory seeds.
\item
  \textbf{Sample size.} Precision-based. An exploratory pilot (draft
  parameters, seeds 900002 to 900011) gave a paired-difference SD of
  0.01587 rad and a projected 95 \% half-width of 0.0114 rad at 10
  seeds, below SESOI/2 = 0.025 rad, so 10 seeds were kept.
\item
  \textbf{Analysis.} 95 \% percentile bootstrap over seeds (10,000
  resamples, bootstrap seed 7) of the paired mean difference, the paired
  t interval, and d\_z. SESOI 0.05 rad. Decision rule: ``reject H0'' if
  the CI excludes 0; ``statistically detectable, practically
  negligible'' if it excludes 0 and lies inside {[}-0.05, 0.05{]};
  ``practically equivalent'' if it includes 0 and lies inside the band;
  otherwise ``inconclusive''. Secondary drift contrasts (POP, MMR,
  EXPLORE and SIM versus RAND) get Holm-adjusted intervals; other
  metrics are descriptive. The Holm construction and the two readings of
  ``practically equivalent'' (band and label) were fixed before any
  confirmatory output (DEC-012) and ratified by the owner (DEC-014); the
  two readings agreed in every verdict.
\item
  \textbf{Pre-registered rules for EXP-003 and EXP-004.} EXP-003: the
  EXP-002 interpretation is retained if every cell with eta
  \textgreater= 0.05 has the EXP-002 sign with a CI excluding 0;
  withdrawn as assumption-dependent if any cell at eta = 0.05 has the
  opposite sign with a CI excluding 0 or is practically equivalent to 0;
  partially robust otherwise. A second check compares oracle and
  estimated profiles. EXP-004: replicates if Delta\_FMA has the EXP-002
  sign with a CI excluding 0; does not replicate if its CI excludes 0
  with the opposite sign or it is practically equivalent to 0; otherwise
  inconclusive.
\item
  \textbf{Exclusions and stopping.} No listener or seed is excluded;
  runs that fail are reported; more than 10 \% failures in an arm stops
  the experiment. Fixed matrices, no optional stopping.
\item
  \textbf{Registration and freeze.} The owner froze protocol v1.0 on
  2026-11-24 (DEC-010; \texttt{research/freeze.json}, revision 17ac055;
  SHA-256 of \texttt{research/protocol.md} a81ed32c\ldots f8c5dc). The
  run guard refuses any confirmatory config whose hash differs from the
  freeze record. SIM minus RAND values seen before the freeze
  (exploratory seeds, draft parameters) are disclosed: pilot mean
  -0.15881 rad (protocol 10.2 item 6). No simulation on the real
  catalogue was run before the freeze.
\item
  \textbf{Execution.} All 560 confirmatory runs were made at code
  revision afc0850 with a clean tree, on 2026-11-24 23:39 to 23:52 UTC,
  on an Apple M1 Pro shared with other sessions (runtimes are upper
  bounds). 560 of 560 runs completed; none failed, was retried or was
  excluded.
\item
  \textbf{Exploratory work, labelled.} EXP-001 and its FMA-path smoke;
  and an M4 confounder inspection (journal J-0008,
  \texttt{results/exploratory/M4-confounders/}) that reads committed
  results and regenerates only the input catalogues and initial tastes.
  It runs no simulation and changes no verdict.
\item
  \textbf{Resources.} Zero cost; local CPU; wall-clock 4 min 51 s
  (EXP-002), 12 min 45 s (EXP-003), 4 min 18 s (EXP-004).
\end{itemize}

\subsection{5. Results}\label{results}

Every number in this section is generated from
\texttt{results/EXP-00N/analysis/analysis.json} (confirmatory) or
\texttt{results/exploratory/M4-confounders/confounders.json}
(exploratory, labelled). The mapping is in
\texttt{research/publication/artifact\_manifest.json}.

\subsubsection{5.1 EXP-002: the primary contrast on the calibrated
synthetic
catalogue}\label{exp-002-the-primary-contrast-on-the-calibrated-synthetic-catalogue}

\textbf{Negative control.} With fixed preferences, drift was exactly 0
for every listener in all 50 fixed runs (claim C-003). This checks the
implementation, not the model.

\textbf{Primary.} Delta (SIM minus RAND) = -0.0355 rad, SD of paired
differences 0.0286, 95 \% bootstrap CI {[}-0.0535, -0.0200{]}, paired t
CI {[}-0.0560, -0.0151{]}, d\_z -1.24, n = 10 seeds. Frozen decision:
\textbf{reject H0}. All 10 seeds had a negative difference (exploratory
sign count, J-0008). The point estimate is inside the SESOI, but the CI
extends past -0.05 rad, so the effect could be smaller or larger than
the smallest effect of interest: its practical size is unresolved.

\textbf{Table 1.} Drift contrasts versus RAND, adaptive, seed-paired
(rad, n = 10 seeds). Source:
\texttt{results/EXP-002/analysis/contrasts.csv}; figure
\texttt{results/EXP-002/analysis/fig-contrasts.svg}.

\begin{longtable}[]{@{}
  >{\raggedright\arraybackslash}p{(\linewidth - 10\tabcolsep) * \real{0.1667}}
  >{\raggedright\arraybackslash}p{(\linewidth - 10\tabcolsep) * \real{0.1667}}
  >{\raggedright\arraybackslash}p{(\linewidth - 10\tabcolsep) * \real{0.1667}}
  >{\raggedright\arraybackslash}p{(\linewidth - 10\tabcolsep) * \real{0.1667}}
  >{\raggedright\arraybackslash}p{(\linewidth - 10\tabcolsep) * \real{0.1667}}
  >{\raggedright\arraybackslash}p{(\linewidth - 10\tabcolsep) * \real{0.1667}}@{}}
\toprule\noalign{}
\begin{minipage}[b]{\linewidth}\raggedright
Contrast
\end{minipage} & \begin{minipage}[b]{\linewidth}\raggedright
Mean diff
\end{minipage} & \begin{minipage}[b]{\linewidth}\raggedright
95 \% bootstrap CI
\end{minipage} & \begin{minipage}[b]{\linewidth}\raggedright
Holm CI
\end{minipage} & \begin{minipage}[b]{\linewidth}\raggedright
d\_z
\end{minipage} & \begin{minipage}[b]{\linewidth}\raggedright
Decision
\end{minipage} \\
\midrule\noalign{}
\endhead
\bottomrule\noalign{}
\endlastfoot
SIM-RAND (primary) & -0.0355 & {[}-0.0535, -0.0200{]} & {[}-0.0535,
-0.0200{]} & -1.24 & reject H0 \\
POP-RAND & 0.1630 & {[}0.1443, 0.1810{]} & {[}0.1418, 0.1834{]} & 5.22 &
reject H0 \\
MMR-RAND & -0.1162 & {[}-0.1339, -0.1009{]} & {[}-0.1377, -0.0983{]} &
-4.18 & reject H0 \\
EXPLORE-RAND & -0.0381 & {[}-0.0570, -0.0219{]} & {[}-0.0625, -0.0186{]}
& -1.27 & reject H0 \\
\end{longtable}

\textbf{Table 2.} Per-policy seed-level means, adaptive preferences
(drift with 95 \% bootstrap CI). Source:
\texttt{results/EXP-002/analysis/arms.csv}; figure
\texttt{results/EXP-002/analysis/fig-arms-drift.svg}.

\begin{longtable}[]{@{}
  >{\raggedright\arraybackslash}p{(\linewidth - 10\tabcolsep) * \real{0.1667}}
  >{\raggedright\arraybackslash}p{(\linewidth - 10\tabcolsep) * \real{0.1667}}
  >{\raggedright\arraybackslash}p{(\linewidth - 10\tabcolsep) * \real{0.1667}}
  >{\raggedright\arraybackslash}p{(\linewidth - 10\tabcolsep) * \real{0.1667}}
  >{\raggedright\arraybackslash}p{(\linewidth - 10\tabcolsep) * \real{0.1667}}
  >{\raggedright\arraybackslash}p{(\linewidth - 10\tabcolsep) * \real{0.1667}}@{}}
\toprule\noalign{}
\begin{minipage}[b]{\linewidth}\raggedright
Policy
\end{minipage} & \begin{minipage}[b]{\linewidth}\raggedright
Drift (rad) {[}95 \% CI{]}
\end{minipage} & \begin{minipage}[b]{\linewidth}\raggedright
Homogenisation
\end{minipage} & \begin{minipage}[b]{\linewidth}\raggedright
Discovery
\end{minipage} & \begin{minipage}[b]{\linewidth}\raggedright
Skip rate
\end{minipage} & \begin{minipage}[b]{\linewidth}\raggedright
Exposure Gini
\end{minipage} \\
\midrule\noalign{}
\endhead
\bottomrule\noalign{}
\endlastfoot
RAND & 1.1119 {[}1.0805, 1.1404{]} & 0.358 & 0.376 & 0.006 & 0.040 \\
POP & 1.2749 {[}1.2513, 1.2966{]} & 0.901 & 0.664 & 0.052 & 0.950 \\
SIM & 1.0763 {[}1.0467, 1.1035{]} & 0.244 & 0.185 & 0.001 & 0.552 \\
MMR & 0.9956 {[}0.9710, 1.0181{]} & 0.219 & 0.290 & 0.002 & 0.378 \\
EXPLORE & 1.0738 {[}1.0446, 1.1005{]} & 0.243 & 0.190 & 0.001 & 0.440 \\
\end{longtable}

\textbf{Exposure (descriptive secondaries, seed-paired SIM minus RAND,
adaptive).} Exposure Gini +0.5118 {[}0.4883, 0.5336{]}; discovery
-0.1909 {[}-0.2119, -0.1698{]}; intra-list diversity -0.6727 {[}-0.6948,
-0.6512{]}; homogenisation -0.1145 {[}-0.1572, -0.0754{]}. These are
descriptive, not confirmatory.

\textbf{Uniform-random exposure is not neutral.} Under the adaptive
model, RAND itself moved preferences by 1.1119 rad on average, closer to
the angle between unrelated directions than to 0. The question the
primary contrast answers is therefore ``which policy moves taste more'',
not ``which policy moves taste at all''.

\subsubsection{5.2 EXP-003: sensitivity to the update rule, rate and
profile}\label{exp-003-sensitivity-to-the-update-rule-rate-and-profile}

All 18 cells are in Table S2 (supplementary) and
\texttt{results/EXP-003/analysis/fig-contrasts.svg}. The cell consumed /
0.05 / estimated has the EXP-002 configuration and reproduces its
primary contrast exactly (-0.0355 rad, CI {[}-0.0535, -0.0200{]}).

\textbf{Table 3.} SIM minus RAND drift by update rule and rate,
estimated profile (rad, 95 \% bootstrap CI, n = 10 seeds). Source:
\texttt{results/EXP-003/analysis/contrasts.csv}.

\begin{longtable}[]{@{}
  >{\raggedright\arraybackslash}p{(\linewidth - 6\tabcolsep) * \real{0.2500}}
  >{\raggedright\arraybackslash}p{(\linewidth - 6\tabcolsep) * \real{0.2500}}
  >{\raggedright\arraybackslash}p{(\linewidth - 6\tabcolsep) * \real{0.2500}}
  >{\raggedright\arraybackslash}p{(\linewidth - 6\tabcolsep) * \real{0.2500}}@{}}
\toprule\noalign{}
\begin{minipage}[b]{\linewidth}\raggedright
Rule
\end{minipage} & \begin{minipage}[b]{\linewidth}\raggedright
eta 0.01
\end{minipage} & \begin{minipage}[b]{\linewidth}\raggedright
eta 0.05
\end{minipage} & \begin{minipage}[b]{\linewidth}\raggedright
eta 0.2
\end{minipage} \\
\midrule\noalign{}
\endhead
\bottomrule\noalign{}
\endlastfoot
consumed & 0.1364 {[}0.1234, 0.1484{]} & -0.0355 {[}-0.0535, -0.0200{]}
& -0.0978 {[}-0.1186, -0.0770{]} \\
exposure & 0.1710 {[}0.1421, 0.1983{]} & -0.1510 {[}-0.1838, -0.1257{]}
& -0.2688 {[}-0.3046, -0.2354{]} \\
bounded & 0.0363 {[}0.0275, 0.0466{]} & 0.0069 {[}-0.0130, 0.0303{]} &
-0.1012 {[}-0.1318, -0.0617{]} \\
\end{longtable}

\textbf{Verdict under the frozen rule: withdrawn as
assumption-dependent} (both equivalence readings agree; claim C-004).
Two cells at eta 0.05 trigger withdrawal: bounded / estimated is
practically equivalent to 0 (+0.0069 rad, CI {[}-0.0130, 0.0303{]}; the
paired t interval {[}-0.0194, 0.0331{]} also lies inside the band), and
bounded / oracle has the opposite sign (+0.1883 rad, CI {[}0.1598,
0.2196{]}; Holm CI {[}0.1478, 0.2362{]}). Beyond the trigger, the sign
changes with eta under the consumed and exposure rules: positive at
0.01, negative at 0.05 and 0.2, with all 10 seeds sharing the sign in
each of these cells (J-0008).

\textbf{Oracle versus estimated profile} (protocol section 7, second
check; C-005). The point signs differ, with both CIs excluding 0, in one
rule-rate pair: bounded, eta 0.2 (estimated -0.1012 rad, oracle +0.1564
rad). There, the effect comes from the profile estimator rather than the
policy family. In the other eight pairs the point signs agree.

\subsubsection{5.3 EXP-004: second test on the real FMA
catalogue}\label{exp-004-second-test-on-the-real-fma-catalogue}

\textbf{Negative control.} Drift was exactly 0 in all 50 fixed runs
(C-003).

\textbf{Primary.} Delta\_FMA = +0.1652 rad, SD 0.0148, 95 \% CI
{[}0.1560, 0.1735{]}, paired t CI {[}0.1546, 0.1758{]}, d\_z 11.15, n =
10 seeds (all 10 positive). Frozen decision: reject H0, with the
opposite sign to EXP-002. \textbf{Frozen reading: does not replicate}
(both readings agree; C-006).

\textbf{Table 4.} Drift contrasts versus RAND on the FMA catalogue (rad,
n = 10 seeds). Source: \texttt{results/EXP-004/analysis/contrasts.csv};
figure \texttt{results/EXP-004/analysis/fig-contrasts.svg}.

\begin{longtable}[]{@{}
  >{\raggedright\arraybackslash}p{(\linewidth - 10\tabcolsep) * \real{0.1667}}
  >{\raggedright\arraybackslash}p{(\linewidth - 10\tabcolsep) * \real{0.1667}}
  >{\raggedright\arraybackslash}p{(\linewidth - 10\tabcolsep) * \real{0.1667}}
  >{\raggedright\arraybackslash}p{(\linewidth - 10\tabcolsep) * \real{0.1667}}
  >{\raggedright\arraybackslash}p{(\linewidth - 10\tabcolsep) * \real{0.1667}}
  >{\raggedright\arraybackslash}p{(\linewidth - 10\tabcolsep) * \real{0.1667}}@{}}
\toprule\noalign{}
\begin{minipage}[b]{\linewidth}\raggedright
Contrast
\end{minipage} & \begin{minipage}[b]{\linewidth}\raggedright
Mean diff
\end{minipage} & \begin{minipage}[b]{\linewidth}\raggedright
95 \% bootstrap CI
\end{minipage} & \begin{minipage}[b]{\linewidth}\raggedright
Holm CI
\end{minipage} & \begin{minipage}[b]{\linewidth}\raggedright
d\_z
\end{minipage} & \begin{minipage}[b]{\linewidth}\raggedright
Decision
\end{minipage} \\
\midrule\noalign{}
\endhead
\bottomrule\noalign{}
\endlastfoot
SIM-RAND (primary) & 0.1652 & {[}0.1560, 0.1735{]} & {[}0.1547,
0.1749{]} & 11.15 & reject H0 \\
POP-RAND & -0.0357 & {[}-0.0723, 0.0005{]} & {[}-0.0723, 0.0005{]} &
-0.58 & inconclusive \\
MMR-RAND & 0.0570 & {[}0.0502, 0.0643{]} & {[}0.0490, 0.0660{]} & 4.76 &
reject H0 \\
EXPLORE-RAND & 0.1631 & {[}0.1539, 0.1720{]} & {[}0.1516, 0.1744{]} &
10.66 & reject H0 \\
\end{longtable}

\textbf{Table 5.} Per-policy seed-level means on FMA, adaptive. Source:
\texttt{results/EXP-004/analysis/arms.csv}; figure
\texttt{results/EXP-004/analysis/fig-arms-drift.svg}.

\begin{longtable}[]{@{}
  >{\raggedright\arraybackslash}p{(\linewidth - 10\tabcolsep) * \real{0.1667}}
  >{\raggedright\arraybackslash}p{(\linewidth - 10\tabcolsep) * \real{0.1667}}
  >{\raggedright\arraybackslash}p{(\linewidth - 10\tabcolsep) * \real{0.1667}}
  >{\raggedright\arraybackslash}p{(\linewidth - 10\tabcolsep) * \real{0.1667}}
  >{\raggedright\arraybackslash}p{(\linewidth - 10\tabcolsep) * \real{0.1667}}
  >{\raggedright\arraybackslash}p{(\linewidth - 10\tabcolsep) * \real{0.1667}}@{}}
\toprule\noalign{}
\begin{minipage}[b]{\linewidth}\raggedright
Policy
\end{minipage} & \begin{minipage}[b]{\linewidth}\raggedright
Drift (rad) {[}95 \% CI{]}
\end{minipage} & \begin{minipage}[b]{\linewidth}\raggedright
Homogenisation
\end{minipage} & \begin{minipage}[b]{\linewidth}\raggedright
Discovery
\end{minipage} & \begin{minipage}[b]{\linewidth}\raggedright
Skip rate
\end{minipage} & \begin{minipage}[b]{\linewidth}\raggedright
Exposure Gini
\end{minipage} \\
\midrule\noalign{}
\endhead
\bottomrule\noalign{}
\endlastfoot
RAND & 0.6933 {[}0.6794, 0.7070{]} & 0.159 & 0.713 & 0.011 & 0.047 \\
POP & 0.6576 {[}0.6205, 0.6893{]} & 0.503 & 0.694 & 0.325 & 0.961 \\
SIM & 0.8585 {[}0.8427, 0.8739{]} & 0.187 & 0.660 & 0.001 & 0.726 \\
MMR & 0.7503 {[}0.7409, 0.7606{]} & 0.146 & 0.656 & 0.001 & 0.477 \\
EXPLORE & 0.8564 {[}0.8395, 0.8730{]} & 0.180 & 0.658 & 0.001 & 0.579 \\
\end{longtable}

\textbf{Exposure (descriptive, SIM minus RAND, adaptive).} Exposure Gini
+0.6796 {[}0.6652, 0.6925{]}; discovery -0.0539 {[}-0.0657, -0.0422{]};
intra-list diversity -0.5309 {[}-0.5429, -0.5198{]}; homogenisation
+0.0275 {[}0.0118, 0.0455{]}, the opposite sign to EXP-002.

\textbf{POP.} Its drift contrast is inconclusive, but POP's skip rate is
0.325 against 0.011 for RAND. Per accepted item, POP drifts more than
RAND (+0.002818 rad per track, CI {[}0.002458, 0.003111{]}; exploratory
side-by-side, J-0008): the sign of the POP contrast depends on whether
skips are counted, so the skip-rate confounder dominates it.

\subsubsection{5.4 Robustness and confounder inspection (exploratory,
J-0008)}\label{robustness-and-confounder-inspection-exploratory-j-0008}

\begin{itemize}
\tightlist
\item
  \textbf{Skip rate.} For SIM, MMR and EXPLORE, drift and drift per
  accepted item have the same sign in both catalogues, and skip-rate
  differences are at most about 0.01. Skips do not explain the SIM
  contrasts. They do explain the POP contrast on FMA (section 5.3).
\item
  \textbf{Seed consistency.} In EXP-002 and EXP-004 all 10 seeds share
  the sign of the SIM, MMR and EXPLORE contrasts. The
  withdrawal-triggering cell bounded / 0.05 / estimated is split (6
  negative, 4 positive), as expected for a near-zero effect.
\item
  \textbf{Ceiling.} Listener-level drift is not capped at pi/2: in
  EXP-002 the median under RAND is 1.2292 rad and 4.2 \% of listener
  values exceed pi/2 (maximum 1.8046 rad). The concern that contrasts
  are compressed by a ceiling (EXP-002 protocol item 16) is weaker than
  the M3 report implied; the M3 phrase ``pi/2 ceiling'' was imprecise
  and is corrected here.
\item
  \textbf{Catalogue geometry (inputs only).} The two catalogues differ
  sharply. In the synthetic catalogue, tracks sit tightly around their
  cluster centres (mean cosine to own centre 0.9313) and the catalogue
  has a strong common direction (resultant length of the mean embedding
  0.4138), because the fitted cluster weights put about 28 \% of tracks
  in one cluster. On FMA, tracks are far from their genre centres
  (0.3327) and spread almost isotropically (resultant length 0.0777),
  while genre centres are closer to one another (mean pairwise cosine
  0.1692 against 0.0082). This is consistent with, but does not
  establish, the reading that RAND drift on the synthetic catalogue is
  driven by the catalogue's common direction, and that SIM on FMA
  follows local neighbourhoods that lead away from the genre-centre
  start points. The design cannot separate embedding, genre balance and
  popularity tail (EXP-004 protocol item 3).
\end{itemize}

\subsubsection{5.5 Expected versus
observed}\label{expected-versus-observed}

\textbf{Table 6.} Predictions, observations and verdicts under the
frozen rules.

\begin{longtable}[]{@{}
  >{\raggedright\arraybackslash}p{(\linewidth - 6\tabcolsep) * \real{0.2500}}
  >{\raggedright\arraybackslash}p{(\linewidth - 6\tabcolsep) * \real{0.2500}}
  >{\raggedright\arraybackslash}p{(\linewidth - 6\tabcolsep) * \real{0.2500}}
  >{\raggedright\arraybackslash}p{(\linewidth - 6\tabcolsep) * \real{0.2500}}@{}}
\toprule\noalign{}
\begin{minipage}[b]{\linewidth}\raggedright
Experiment
\end{minipage} & \begin{minipage}[b]{\linewidth}\raggedright
Prediction or rule (frozen)
\end{minipage} & \begin{minipage}[b]{\linewidth}\raggedright
Observation
\end{minipage} & \begin{minipage}[b]{\linewidth}\raggedright
Verdict
\end{minipage} \\
\midrule\noalign{}
\endhead
\bottomrule\noalign{}
\endlastfoot
EXP-001 & V1 to V6 pass; pilot half-width below 0.025 rad & All checks
pass in two runs, byte-identical across revisions; half-width 0.0114 rad
& Pass (exploratory; no claim) \\
EXP-002 negative control & Drift exactly 0 under fixed preferences & 0
in all 50 fixed runs & Pass \\
EXP-002 primary & H1: Delta != 0 (no sign predicted); disconfirmed if CI
inside +/-0.05 & -0.0355 rad {[}-0.0535, -0.0200{]} & Reject H0;
practical size unresolved \\
EXP-002 secondaries & Holm family of four drift contrasts & POP +, MMR
-, EXPLORE - vs RAND, all Holm CIs exclude 0 & All rejected after
Holm \\
EXP-003 & Retain if every cell with eta \textgreater= 0.05 keeps the
sign with CI excluding 0 & Bounded / 0.05: equivalent (estimated) and
reversed (oracle); signs change with eta & Withdrawn as
assumption-dependent \\
EXP-003 second check & Oracle and estimated signs agree & Differ, both
CIs excluding 0, at bounded / 0.2 & Profile-dependent in one pair \\
EXP-004 primary & Replicates if same sign with CI excluding 0 & +0.1652
rad {[}0.1560, 0.1735{]} & Does not replicate \\
EXP-004 secondaries & Holm family of four & MMR +, EXPLORE + rejected;
POP inconclusive & Signs reversed relative to EXP-002 for MMR and
EXPLORE \\
\end{longtable}

The pre-freeze pilot (draft parameters) had shown SIM below RAND by
-0.15881 rad. The calibrated confirmatory estimate kept the sign but was
much smaller, and its paired SD (0.0286 rad) was larger than the pilot's
(0.01587 rad). The 10-seed design was sized on the pilot SD; the
achieved CI half-width was about 0.017 rad, still within the SESOI/2
target.

\subsection{6. Discussion}\label{discussion}

\textbf{What the evidence supports.} Within this simulation, two
statements survive every test the protocol set. First, recommendation
policies change exposure: SIM concentrates impressions (exposure Gini)
and lowers within-slate diversity relative to RAND on both catalogues,
with consistent signs. This is largely by construction for a
personalised top-k policy, and these contrasts were descriptive
secondaries, so it is a bounded, unsurprising result. Second, when
preferences adapt to consumption, any exposure, including uniform-random
exposure, changes simulated taste, and the policy changes how much.

\textbf{What the evidence does not support.} It does not support a
policy-level statement about the direction of the taste effect. The
primary contrast was negative on the calibrated synthetic catalogue,
reversed on the real catalogue, reversed at a lower update rate, and
vanished or reversed under bounded-confidence updating. The
pre-registered rules turned these into a withdrawal (EXP-003) and a
non-replication (EXP-004). The direction is a joint property of the
policy, the update rule, the update rate and the catalogue geometry.

\textbf{Competing explanations, and which the results address.}

\begin{longtable}[]{@{}
  >{\raggedright\arraybackslash}p{(\linewidth - 6\tabcolsep) * \real{0.2500}}
  >{\raggedright\arraybackslash}p{(\linewidth - 6\tabcolsep) * \real{0.2500}}
  >{\raggedright\arraybackslash}p{(\linewidth - 6\tabcolsep) * \real{0.2500}}
  >{\raggedright\arraybackslash}p{(\linewidth - 6\tabcolsep) * \real{0.2500}}@{}}
\toprule\noalign{}
\begin{minipage}[b]{\linewidth}\raggedright
Conclusion
\end{minipage} & \begin{minipage}[b]{\linewidth}\raggedright
Competing explanation (protocol)
\end{minipage} & \begin{minipage}[b]{\linewidth}\raggedright
Addressed by
\end{minipage} & \begin{minipage}[b]{\linewidth}\raggedright
Status
\end{minipage} \\
\midrule\noalign{}
\endhead
\bottomrule\noalign{}
\endlastfoot
EXP-002 primary & Skip rate (relevance) rather than steering & Skip
rates 0.001 and 0.006; per-accept contrast same sign & Addressed: not
the explanation for SIM \\
EXP-002 primary & Profile estimator rather than policy family & EXP-003
oracle cells: same sign as estimated at consumed / 0.05 (-0.0607 rad) &
Addressed for the primary cell; not for bounded / 0.2 \\
EXP-002 primary & Update rule and rate & EXP-003 & Confirmed as an
explanation: interpretation withdrawn \\
EXP-002 primary & Compression near pi/2 & Listener-level distribution
(J-0008) & Partly addressed: no hard ceiling observed \\
EXP-002 vs EXP-004 & Embedding, genre balance, popularity tail & Not
separable in this design; geometry inspection only & Open \\
EXP-004 POP & Skip rate & Per-accept contrast reverses sign & Confirmed
as an explanation for POP \\
\end{longtable}

\textbf{Relation to prior work.} Homogenisation under feedback loops
{[}chaney2018{]} and popularity amplification {[}mansoury2020{]} are
reported in prior simulations. In ECHO, POP produced the largest
homogenisation on both catalogues (POP minus RAND +0.5433 on synthetic,
+0.3439 on FMA, descriptive), which is in line with those reports, while
SIM's homogenisation effect changed sign between catalogues. Their
models differ from ECHO's (fixed intrinsic preferences, rating-based
simulation), so no finer comparison is made. The result that the sign of
the drift contrast depends on the assumed update rule is consistent with
the motivation in carroll2022 for measuring induced shift against a
reference, and it argues for reporting assumption sensitivity alongside
any single-mechanism simulation.

\textbf{An alternative reading of the drift metric (second review
pass).} Under the consumed-item rule each accepted track multiplies the
weight of the previous preference by about (1 - eta) before
normalisation. With eta = 0.05 and about 100 accepted tracks per
session, the weight left on the initial taste is of the order of
0.95\^{}100, about 0.006. End-of-session ``taste'' is then close to a
recency-weighted average of the last few dozen consumed tracks, so the
drift outcome largely re-measures exposure through consumption rather
than an independent property of the listener. At eta = 0.01 the
corresponding weight is about 0.37, so the initial taste still
dominates, which is one candidate explanation for the sign change with
the rate in EXP-003. This reading is analytical and was not tested; it
means that, in this model, ``exposure'' and ``taste'' outcomes are less
separable than the research question suggests.

\textbf{Practical significance.} The drift unit is an angle in a
16-dimensional synthetic taste space; 0.05 rad (about 2.9 degrees) was a
judgement-based SESOI. No mapping from these angles to human listening
outcomes exists in this study.

\subsection{7. Limitations and threats to
validity}\label{limitations-and-threats-to-validity}

\begin{itemize}
\tightlist
\item
  \textbf{Simulation only.} Listeners are synthetic agents with assumed
  update rules. Nothing here is evidence about human persuasion, real
  listeners or any commercial or proprietary recommender.
\item
  \textbf{Assumed mechanisms.} Three update rules and three rates are a
  small slice of plausible mechanisms. The withdrawal shows that this
  slice already matters.
\item
  \textbf{Catalogue.} The synthetic catalogue is stylised (Gaussian
  clusters with FMA-fitted weights and popularity). The real catalogue
  is one embedding (PCA-16 of librosa statistics) of freely licensed
  music; audio-feature similarity is not perceived similarity. The 10
  FMA samples of 1000 tracks overlap within 4,951 tracks.
\item
  \textbf{Statistics.} 10 seeds per arm; the percentile bootstrap with n
  = 10 can under-cover (the paired t intervals agree in sign). EXP-003's
  withdrawal rule has no multiplicity correction by design; the
  triggering oracle cell also excludes 0 after Holm. Several bootstrap
  p-values are exactly 0, and Holm ties are ranked by name, which
  changes interval levels only. Large d\_z values (11.15 on FMA) reflect
  small between-seed SDs, not large practical effects. On FMA the ten
  1000-track samples are drawn from the same 4,951 tracks, so seeds are
  not fully independent and the EXP-004 intervals may be too narrow; the
  sign of the EXP-004 contrast, shared by all 10 seeds, is less exposed
  to this than its width.
\item
  \textbf{Construct.} Angular drift measures how far taste moved, not
  toward what. The SESOI of 0.05 rad was proposed by the AI and kept by
  the owner; it is not empirically anchored.
\item
  \textbf{Model choices.} \texttt{skip\_utility\ =\ 0} makes skips rare
  except under POP. The implementation choices of DEC-006 (spreads 0.1,
  warm-up for the oracle profile, synchronous POP counts) may shape
  results.
\item
  \textbf{Pre-freeze exposure.} The pilot SIM minus RAND value at draft
  parameters was seen before the freeze and disclosed; the calibrated
  parameters were fitted on FMA metadata without any simulation. The
  statement that the SIM minus RAND values of the exploratory FMA-path
  smoke were not inspected before the confirmatory runs rests on the
  session record and cannot be verified independently; the smoke used
  smoke-size parameters and exploratory seeds.
\item
  \textbf{Construct overlap.} With eta 0.05 the initial taste is almost
  forgotten after 100 updates, so drift largely reflects recent
  consumption (section 6).
\item
  \textbf{Runtime figures} come from a shared, loaded machine and are
  upper bounds.
\item
  \textbf{Literature.} Eight sources, all read in full; non-systematic
  search; citation chaining from kalimeris2021 and carroll2022 was not
  done, so closely related models may have been missed.
\end{itemize}

The full list is \texttt{research/limitations.md}.

\subsection{8. Ethics and responsible
research}\label{ethics-and-responsible-research}

The study uses public FMA metadata under CC BY 4.0, with attribution; no
audio is downloaded, no artist or title fields are processed, and no
personal or listener data exists. No human participants were involved
and no ethics approval was sought or claimed; a study of human listeners
would need an owner-approved protocol with appropriate review. The study
evaluates exposure concentration and taste change as outcomes to be
understood, not targets to optimise; it should not be read as a recipe
for steering preferences. No Spotify or other service was scraped, and
no claim is made about any proprietary system.

\subsection{9. Conclusion and future
work}\label{conclusion-and-future-work}

Within a pre-registered simulation, a similarity recommender changed
exposure relative to uniform-random recommendation on both a calibrated
synthetic and a real music catalogue (descriptive outcomes, largely
expected by construction). Its effect on the amount of simulated taste
change was statistically detectable in the primary experiment, of
unresolved practical size, and its direction was not robust: the frozen
sensitivity rule withdrew it, and the real-catalogue second test
reversed it. The bounded answer to the title question is that, in this
model, whether a recommender ``shapes'' taste more or less than random
exposure depends on how simulated listeners update and on the catalogue,
not on the policy alone.

Future work, each as a new registered experiment: locate the update-rate
threshold at which the SIM minus RAND sign changes for each rule;
separate embedding, genre balance and popularity tail on the real
catalogue (for example by re-embedding FMA tracks or reweighting
genres); add a directional outcome (movement toward or away from the
home cluster); complete citation chaining and full-text reading of the
closest prior models.

\subsection{10. Data, code and artifact
availability}\label{data-code-and-artifact-availability}

\begin{itemize}
\tightlist
\item
  Repository:
  https://github.com/mohamedbechirkefi/echo-recommender-feedback.
\item
  Interactive page:
  https://echo-recommender-feedback.kefimohamedbechir.workers.dev.
\item
  Protocols and freeze record: \texttt{research/protocol.md},
  \texttt{research/experiments/EXP-00N/protocol.md},
  \texttt{research/freeze.json}.
\item
  Experiment reports: \texttt{research/experiments/EXP-00N/report.md}.
\item
  Raw results: \texttt{results/EXP-00N/runs/} (schema
  \texttt{echo.result.v1}), manifests, \texttt{provenance.csv},
  \texttt{SHA256SUMS}.
\item
  Analysis:
  \texttt{PYTHONPATH=src\ .venv/bin/python\ -m\ echo.analysis\ results/EXP-00N\ -\/-planned\ N\ -\/-kind\ \{policy\textbar{}sensitivity\textbar{}replication\}};
  re-running it on the committed runs reproduced every analysis file
  byte for byte on 2026-11-25 (evidence audit). After the owner-approved
  formatting fix of DEC-016, the regenerated tables hold the same
  values.
\item
  Exploratory confounder inspection:
  \texttt{PYTHONPATH=src\ .venv/bin/python\ scripts/m4\_exploratory\_confounders.py}.
\item
  Data: \texttt{scripts/download\_fma\_metadata.py} and
  \texttt{scripts/build\_fma\_catalogue.py} rebuild the git-ignored FMA
  files; the processed catalogue's SHA-256 is checked at run time. Only
  aggregate calibration values are committed.
\item
  Evidence trail: \texttt{research/claims\_register.csv},
  \texttt{research/publication/artifact\_manifest.json},
  \texttt{research/publication/claims\_audit.md}.
\end{itemize}

\subsection{11. Funding, competing interests and AI
assistance}\label{funding-competing-interests-and-ai-assistance}

\begin{itemize}
\tightlist
\item
  \textbf{Funding:} No funding was received. The only expense was the
  owner's existing AI coding subscription.
\item
  \textbf{Competing interests:} The author declares no competing
  interests.
\item
  \textbf{AI assistance:} an AI system (Claude, Anthropic) wrote the
  simulation and analysis code, ran the experiments under the
  owner-frozen protocol and drafted the reports and this manuscript. The
  owner set the question, approved the protocol freeze, the real-data
  use and every post-freeze decision (DEC-010, DEC-014), and retains
  authority over claims and publication. AI review of the manuscript
  (\texttt{research/publication/review\_pass\_2.md}) is internal quality
  control, not peer review. Every number was checked mechanically
  against committed result files (section 10).
\end{itemize}

\subsection{References}\label{references}

Abstract pages were first opened on 2026-11-24. The full texts of all
eight sources were read: vargas2011 on 2026-11-25 (M4), the other seven
on 2026-11-25 (M6), from the free arXiv PDFs or, for kalimeris2021, the
PDF linked from the Meta Research publication page.

\begin{itemize}
\tightlist
\item
  {[}carroll2022{]} M. Carroll, A. Dragan, S. Russell, D.
  Hadfield-Menell. Estimating and Penalizing Induced Preference Shifts
  in Recommender Systems. ICML 2022, PMLR 162:2686-2708.
  arXiv:2204.11966. https://arxiv.org/abs/2204.11966 (full text with
  appendices).
\item
  {[}chaney2018{]} A. J. B. Chaney, B. M. Stewart, B. E. Engelhardt. How
  Algorithmic Confounding in Recommendation Systems Increases
  Homogeneity and Decreases Utility. arXiv:1710.11214, 2018.
  https://arxiv.org/abs/1710.11214 (full text, 9 pages; the RecSys 2018
  DOI listed on arXiv was not opened).
\item
  {[}defferrard2017{]} M. Defferrard, K. Benzi, P. Vandergheynst, X.
  Bresson. FMA: A Dataset For Music Analysis. ISMIR 2017;
  arXiv:1612.01840. https://arxiv.org/abs/1612.01840 and
  https://github.com/mdeff/fma (full text, 8 pages, and repository
  README).
\item
  {[}ie2019{]} E. Ie, C. Hsu, M. Mladenov, V. Jain, S. Narvekar, J.
  Wang, R. Wu, C. Boutilier. RecSim: A Configurable Simulation Platform
  for Recommender Systems. arXiv:1909.04847, 2019.
  https://arxiv.org/abs/1909.04847 (full text).
\item
  {[}jiang2019{]} R. Jiang, S. Chiappa, T. Lattimore, A. György, P.
  Kohli. Degenerate Feedback Loops in Recommender Systems. AIES 2019.
  DOI 10.1145/3306618.3314288. https://arxiv.org/abs/1902.10730 (full
  text, 10 pages).
\item
  {[}kalimeris2021{]} D. Kalimeris, S. Bhagat, S. Kalyanaraman, U.
  Weinsberg. Preference Amplification in Recommender Systems. KDD 2021.
  https://research.facebook.com/publications/preference-amplification-in-recommender-systems/
  (full text, 11 pages; DOI 10.1145/3447548.3467298 as printed in the
  paper, not opened).
\item
  {[}mansoury2020{]} M. Mansoury, H. Abdollahpouri, M. Pechenizkiy, B.
  Mobasher, R. Burke. Feedback Loop and Bias Amplification in
  Recommender Systems. arXiv:2007.13019, 2020.
  https://arxiv.org/abs/2007.13019 (full text, 5 pages).
\item
  {[}vargas2011{]} S. Vargas, P. Castells. Rank and Relevance in Novelty
  and Diversity Metrics for Recommender Systems. RecSys 2011, Chicago.
  https://castells.github.io/papers/recsys2011.pdf (full text, 8 pages,
  read 2026-11-25; DOI 10.1145/2043932.2043955 taken from a search
  result, not opened).
\end{itemize}

\subsection{Supplementary appendix}\label{supplementary-appendix}

See \texttt{research/paper/supplementary.md}: the complete experimental
matrix, Table S2 (all 18 EXP-003 cells), the oracle versus estimated
table, the exploratory confounder tables, configurations and hashes,
report links, and every deviation and owner decision.

\end{document}
